Un cuerpo de temperatura uniforme tiene una capacidad térmica desconocida \(C>0\), en \(\mathrm{J/K}\): para aumentar su temperatura en \(1\) K necesita \(C\) julios. Un calentador le aporta una potencia conocida \(P\). El cuerpo pierde hacia un ambiente a \(20\,{}^\circ\mathrm C\) una potencia proporcional a la diferencia de temperaturas, con coeficiente desconocido \(H>0\), en \(\mathrm{W/K}\). No hay otros intercambios.

En dos ensayos simulados del mismo cuerpo se obtienen estos datos de arranque, tratados como exactos:
\[
\begin{array}{c|ccc}
\text{ensayo}&T(0)\ ({}^\circ\mathrm C)&P(0)\ (\mathrm W)&T'(0)\ (\mathrm{K/s})\\\hline
1&20&12&0.30\\
2&30&0&-0.50
\end{array}
\]
Elija una variable de estado, formule el balance y determine \(C\), \(H\) y el tiempo característico de relajación. Considere después dos ensayos generales con potencias y temperaturas iniciales elegidas por el experimentador. Suponiendo que los datos respetan el modelo, establezca qué condición sobre esos arranques permite identificar \(C\) y \(H\) por separado mediante sus pendientes iniciales. Explique qué información se perdería si ambos ensayos comenzaran a la temperatura ambiente.
